\usepackage[utf8]{inputenc}
\usepackage[T1]{fontenc}
\usepackage{lmodern}
\usepackage[polish,provide=*]{babel}
\usepackage[margin=2.5cm]{geometry}
\usepackage{amsmath,amssymb,amsthm,mathtools}
\usepackage{array}
\usepackage{enumitem}
\usepackage{float}
\usepackage{makeidx}
\usepackage{tikz}
\usepackage{tikz-cd}
\usetikzlibrary{arrows.meta,calc,positioning,decorations.pathreplacing,decorations.markings,patterns}
\usepackage{pgfplots}
\pgfplotsset{compat=1.18}
\definecolor{manifoldblue}{RGB}{35,91,140}
\definecolor{manifoldred}{RGB}{170,65,55}
\definecolor{manifoldgreen}{RGB}{55,125,85}
\definecolor{manifoldgold}{RGB}{205,145,35}
\usepackage{fancyhdr}
\usepackage{hyperref}
\hypersetup{colorlinks=true, linkcolor=blue!45!black, citecolor=blue!45!black, urlcolor=blue!45!black}
\makeindex

\newcommand{\R}{\mathbb{R}}
\newcommand{\C}{\mathbb{C}}
\newcommand{\N}{\mathbb{N}}
\newcommand{\Z}{\mathbb{Z}}
\newcommand{\Q}{\mathbb{Q}}
\newcommand{\id}{\mathrm{id}}
\newcommand{\Lie}{\mathcal L}
\newcommand{\X}{\mathfrak X}
\newcommand{\supp}{\operatorname{supp}}
\newcommand{\pr}{\operatorname{pr}}
\newcommand{\D}{\mathrm D}
\newcommand{\dd}{\mathrm d}
\renewcommand{\Re}{\operatorname{Re}}
\renewcommand{\Im}{\operatorname{Im}}

\theoremstyle{plain}
\newtheorem{twierdzenie}{Twierdzenie}[section]
\newtheorem{lemat}[twierdzenie]{Lemat}
\newtheorem{stwierdzenie}[twierdzenie]{Stwierdzenie}
\newtheorem{wniosek}[twierdzenie]{Wniosek}
\newtheorem{fakt}[twierdzenie]{Fakt}
\theoremstyle{definition}
\newtheorem{definicja}[twierdzenie]{Definicja}
\newtheorem{przyklad}[twierdzenie]{Przykład}
\theoremstyle{remark}
\newtheorem{uwaga}[twierdzenie]{Uwaga}
\renewcommand{\proofname}{Dowód}

\pagestyle{fancy}
\fancyhf{}
\fancyhead[L]{\small Geometria i topologia różniczkowa}
\fancyhead[R]{\small\ifnum\value{chapter}>0 Rozdział \thechapter\else Wprowadzenie\fi}
\fancyfoot[C]{\thepage}
\renewcommand{\headrulewidth}{0.4pt}

\setlength{\parindent}{0pt}
\setlength{\parskip}{4pt}
\setlength{\emergencystretch}{2em}
